\documentclass[10pt]{article}
% Shared LaTeX style for MIT 18.06 exam revision notes.
% Compile topic files with Tectonic, XeLaTeX, or LuaLaTeX.

\usepackage[a4paper,margin=0.72in]{geometry}
\usepackage{fontspec}
\usepackage{amsmath,amssymb,mathtools}
\usepackage{booktabs,array,tabularx}
\usepackage{enumitem}
\usepackage{titlesec}
\usepackage{fancyhdr}
\usepackage{microtype}
\usepackage{xcolor}

\setmainfont{Times New Roman}
\setsansfont{Arial}

\setlength{\parindent}{0pt}
\setlength{\parskip}{3.6pt}
\setlength{\abovedisplayskip}{6pt}
\setlength{\belowdisplayskip}{6pt}
\setlength{\abovedisplayshortskip}{4pt}
\setlength{\belowdisplayshortskip}{4pt}
\renewcommand{\arraystretch}{1.08}
\setlength{\tabcolsep}{4.5pt}

\titleformat{\section}
  {\normalfont\large\bfseries}
  {\thesection}
  {0.55em}
  {}
\titlespacing*{\section}{0pt}{10pt}{3pt}

\titleformat{\subsection}
  {\normalfont\normalsize\bfseries}
  {\thesubsection}
  {0.5em}
  {}
\titlespacing*{\subsection}{0pt}{7pt}{2pt}

\setlist[itemize]{leftmargin=1.35em,itemsep=1pt,topsep=2pt,parsep=0pt}
\setlist[enumerate]{leftmargin=1.55em,itemsep=1pt,topsep=2pt,parsep=0pt}

\pagestyle{fancy}
\fancyhf{}
\fancyfoot[C]{\scriptsize\color{black!45}MIT 18.06 Linear Algebra Exam Notes --- \thepage}
\renewcommand{\headrulewidth}{0pt}
\renewcommand{\footrulewidth}{0pt}

\newcommand{\notesubtitle}[1]{%
  {\centering\itshape #1\par}
  \vspace{2pt}
}

\newcommand{\source}[1]{%
  {\centering\small #1\par}
  \vspace{6pt}
}

\newcommand{\labelnote}[2]{%
  \par\smallskip
  \noindent\textbf{#1.}\quad #2
  \par\smallskip
}

\newcommand{\tightlabel}[1]{%
  \par\smallskip
  \noindent\textbf{#1.}\quad
}

\newcolumntype{Y}{>{\raggedright\arraybackslash}X}
\newcolumntype{L}[1]{>{\raggedright\arraybackslash}p{#1}}

\newenvironment{notetable}
  {\small\begin{tabularx}{\linewidth}}
  {\end{tabularx}}


\begin{document}

\begin{center}
{\LARGE 4.2--4.3 Projections and Least Squares\par}
\end{center}
\notesubtitle{Concise exam revision notes for projections, projection matrices, and least squares.}
\source{Based on handwritten MIT 18.06 revision notes.}

\section{Core Idea}

\labelnote{Core idea}{When \(Ax=b\) has no solution, replace \(b\) by the closest vector \(p\) in the column space of \(A\). The error \(e=b-p\) must be perpendicular to that column space.}

The picture is always
\[
b=p+e,
\qquad
p\in C(A),
\qquad
e\in N(A^{T}).
\]
Projection turns an unsolvable equation into a solvable normal equation.

\section{Projection Onto a Line}

For a nonzero vector \(a\), the projection of \(b\) onto the line through \(a\) has the form
\[
p=\hat{x}a.
\]
The error is perpendicular to \(a\):
\[
a^{T}(b-\hat{x}a)=0.
\]
Solving gives
\[
\hat{x}=\frac{a^{T}b}{a^{T}a},
\qquad
p=a\frac{a^{T}b}{a^{T}a}.
\]
So the one-dimensional projection matrix is
\[
P=\frac{aa^{T}}{a^{T}a}.
\]

\labelnote{Meaning}{\(\hat{x}\) is the coordinate along the line; \(p\) is the closest point on the line.}

\section{Projection Onto a Subspace}

Let the columns of \(A\) span the subspace. We look for
\[
p=A\hat{x}.
\]
The error \(e=b-A\hat{x}\) must be perpendicular to every column of \(A\):
\[
A^{T}e=0.
\]
Therefore
\[
A^{T}(b-A\hat{x})=0,
\qquad
A^{T}A\hat{x}=A^{T}b.
\]
These are the normal equations.

If the columns of \(A\) are independent, then \(A^{T}A\) is invertible and
\[
\hat{x}=(A^{T}A)^{-1}A^{T}b.
\]
The projection is
\[
p=A\hat{x}=A(A^{T}A)^{-1}A^{T}b.
\]
So the projection matrix is
\[
P=A(A^{T}A)^{-1}A^{T}.
\]

\section{Projection Matrix Facts}

\begin{center}
\small
\begin{tabularx}{\linewidth}{@{}L{0.25\linewidth}Y Y@{}}
\toprule
Fact & Formula & Meaning \\
\midrule
Projection lands in \(C(A)\) & \(p=Pb\) & \(p\) is the closest vector in the column space. \\
Error is perpendicular & \(A^{T}(b-p)=0\) & \(e=b-p\) lies in \(N(A^{T})\). \\
Projecting twice changes nothing & \(P^{2}=P\) & Once on the subspace, it stays there. \\
Orthogonal projection is symmetric & \(P^{T}=P\) & True for projection onto a column space. \\
\bottomrule
\end{tabularx}
\end{center}

\labelnote{Exam trigger}{If the problem asks for closest vector, best fit, or least squares, expect \(A^{T}A\hat{x}=A^{T}b\).}

\section{Orthogonal Columns Shortcut}

If the columns \(q_1,\ldots,q_n\) are orthonormal, put them in \(Q\). Then
\[
Q^{T}Q=I.
\]
Projection becomes much simpler:
\[
p=QQ^{T}b.
\]
Equivalently,
\[
p=(q_1^{T}b)q_1+\cdots+(q_n^{T}b)q_n.
\]

\begin{center}
\small
\begin{tabularx}{\linewidth}{@{}L{0.28\linewidth}Y@{}}
\toprule
Columns of \(A\) & Projection formula \\
\midrule
Independent & \(P=A(A^{T}A)^{-1}A^{T}\). \\
Orthogonal but not unit & Add projections along each column separately. \\
Orthonormal & \(P=QQ^{T}\). \\
\bottomrule
\end{tabularx}
\end{center}

\labelnote{Common mistake}{Do not use \(QQ^{T}\) unless the columns are orthonormal. Orthogonal columns still need length factors.}

\section{Least Squares}

Least squares solves the closest version of an inconsistent system:
\[
Ax\approx b.
\]
The best coefficient vector is \(\hat{x}\), found from
\[
A^{T}A\hat{x}=A^{T}b.
\]
Then
\[
p=A\hat{x},
\qquad
e=b-p.
\]

For a best-fit line
\[
y=C+Dt,
\]
the matrix has rows
\[
\begin{bmatrix}1&t_i\end{bmatrix}.
\]
The unknown vector is
\[
\hat{x}=
\begin{bmatrix}
C\\
D
\end{bmatrix}.
\]
The same normal equation gives the least-squares line.

\section{\(P\) and \(I-P\)}

The projection splits \(b\) into two perpendicular parts:
\[
b=Pb+(I-P)b.
\]
Here
\[
Pb\in C(A),
\qquad
(I-P)b\in N(A^{T}).
\]
If \(P\) projects onto \(C(A)\), then \(I-P\) projects onto the left nullspace.

\labelnote{Exam memory}{Projection onto a subspace plus projection onto its orthogonal complement gives the original vector.}

\section{Common Mistakes}

\begin{center}
\small
\begin{tabularx}{\linewidth}{@{}L{0.40\linewidth}Y@{}}
\toprule
Mistake & Correct exam habit \\
\midrule
Solving \(Ax=b\) when it is inconsistent. & Solve \(A^{T}A\hat{x}=A^{T}b\). \\
Forgetting the error condition. & Check \(A^{T}e=0\). \\
Using \(P=A^{-1}\) language. & Projection is not inversion. Use \(P=A(A^{T}A)^{-1}A^{T}\). \\
Using \(QQ^{T}\) for non-orthonormal columns. & First check \(Q^{T}Q=I\). \\
Assuming \(P_1+P_2=I\) for any two projections. & This works for orthogonal complementary subspaces. \\
\bottomrule
\end{tabularx}
\end{center}

\section{Final Exam Checklist}

\tightlabel{Checklist}
\begin{enumerate}
  \item Identify the subspace: usually \(C(A)\) or a line through \(a\).
  \item Write \(p=A\hat{x}\) or \(p=\hat{x}a\).
  \item Write the error \(e=b-p\).
  \item Impose perpendicularity: \(A^{T}e=0\) or \(a^{T}e=0\).
  \item Solve \(A^{T}A\hat{x}=A^{T}b\).
  \item Compute \(p=A\hat{x}\).
  \item Compute the error \(e=b-p\).
  \item Use \(P=A(A^{T}A)^{-1}A^{T}\) only when the columns of \(A\) are independent.
  \item Use \(P=QQ^{T}\) only when the columns are orthonormal.
  \item State where \(p\) and \(e\) live: \(p\in C(A)\), \(e\in N(A^{T})\).
\end{enumerate}

\end{document}
